\documentclass[11pt]{article}
\usepackage[a4paper,margin=2.5cm]{geometry}
\usepackage{graphicx}
\usepackage{hyperref}
\usepackage{listings}
\usepackage{xcolor}
\usepackage{booktabs}

\lstset{
  basicstyle=\ttfamily\small,
  breaklines=true,
  frame=single,
  backgroundcolor=\color{gray!10}
}

\title{Cryptography and Network Security \\ \large Integrated Situation --- Technical Report}
\author{Josue Baraka Rudahalingoma \\ ULK Polytechnic Institute}
\date{\today}

\begin{document}
\maketitle

\section{Introduction}
This report documents the security assessment and technical work carried
out for the ETTCS801 Integrated Situation. The scenario involves a
polytechnic institute storing student records on a central server and
transferring files between two campuses, with weak staff passwords,
outdated software, unencrypted transfers, and guest network access to
the records server identified as issues in a security review.

\section{Risk Assessment}
% a. Assets, vulnerabilities, consequences
% b. Risk ranking with reasoning
% c. Recommended controls
% See risk_assessment.md for the full table -- reproduce/summarise it here.

\subsection{Assets, Vulnerabilities and Consequences}
\begin{table}[h]
\centering
\begin{tabular}{lll}
\toprule
Asset & Vulnerability & Consequence \\
\midrule
Records server & Guest network access & Unauthorised access/tampering \\
Staff credentials & Weak passwords & Account compromise \\
File transfer channel & Unencrypted transfer & Interception in transit \\
\bottomrule
\end{tabular}
\caption{Assets, vulnerabilities and consequences}
\end{table}

\subsection{Risk Ranking}
% Insert your likelihood/impact table and reasoning here.

\subsection{Recommended Controls}
% Insert your controls table here.

\section{Encryption and Integrity Toolkit}
The \texttt{crypto\_toolkit.py} programme was developed to encrypt and
decrypt student records and to verify their integrity.

\subsection{Design}
\begin{itemize}
  \item \textbf{Encryption:} symmetric encryption using \texttt{Fernet}
    (AES-128 in CBC mode with HMAC authentication) from the Python
    \texttt{cryptography} library.
  \item \textbf{Integrity:} SHA-256 hashing via \texttt{hashlib}, with
    the hash saved alongside the file and re-checked on demand.
  \item \textbf{Robustness:} all file operations are wrapped in
    \texttt{try/except} blocks so missing files or invalid input do not
    crash the programme.
\end{itemize}

\subsection{Test Evidence}
% Paste terminal output / screenshots here showing:
%  - successful encryption
%  - successful decryption matching the original
%  - hash computed
%  - integrity check passing on an unmodified file
%  - integrity check failing after the file is modified
\begin{lstlisting}
$ python3 crypto_toolkit.py genkey
[OK] Key generated and saved to secret.key
...
\end{lstlisting}

\section{Network Traffic Filtering}
Firewall rules were applied in the authorised laboratory environment
using \texttt{iptables} (see \texttt{firewall\_rules.sh}).

\subsection{Rules Summary}
\begin{itemize}
  \item Block guest subnet $\rightarrow$ records-server port
  \item Allow staff subnet $\rightarrow$ records-server port
  \item Drop all other inbound traffic to that port
\end{itemize}

\subsection{Test Procedure and Results}
% Reproduce the table from filter_tests.md here.
\begin{table}[h]
\centering
\begin{tabular}{llll}
\toprule
Test & Source & Expected & Result \\
\midrule
1 & Staff subnet & Allowed & \\
2 & Guest subnet & Blocked & \\
3 & External/unauthorised & Blocked & \\
\bottomrule
\end{tabular}
\caption{Firewall test results}
\end{table}

\section{Conclusion}
% Summarise findings and the security posture achieved.
The measures implemented -- network segmentation via firewall rules,
authenticated symmetric encryption, and SHA-256 integrity verification
-- address the vulnerabilities identified in the risk assessment.

\vspace{1em}
\noindent\textbf{GitHub repository:} \url{https://github.com/<your-username>/cryptography-network-security-exam}

\end{document}
